\documentclass{article}
\usepackage[utf8]{inputenc}
\usepackage{graphicx}

\title{YAKE + LLM Local Pipeline}
\author{Student}
\date{\today}

\begin{document}
\maketitle

\section{Introducao}
Descreva o problema e a motivacao para combinar YAKE! com LLMs locais.

\section{Metodologia}
Explique o pipeline: pre-processamento, extracao de keywords, geracao do resumo e avaliacao.

\section{Avaliacao}
Inclua cobertura de keywords, ROUGE e exemplos qualitativos.

\section{Conclusao}
Resuma resultados e trabalho futuro.

\end{document}
